%! Bias in Samples and Surveys — Unit 1, Lesson 1.4 — Exit Ticket (Answer Key)
\documentclass[10pt]{article}
\usepackage{saar-article}
\usepackage{saar-key}   % pulls in -boxes; do NOT also load -boxes
\newcommand{\TallMath}[1]{$\displaystyle #1\rule[-1.4em]{0pt}{3.2em}$}

\begin{document}

\pageheader{Unit 1, Lesson 1.4}{Exit Ticket --- Answer Key}
% NAMESTRIP: no \namedateperiod here — mirrors the blank. Teacher-only prose goes
% in the lesson plan as \begin{teachernote}[Exit Ticket], never in a key.

\vspace{0.15in}

\begin{notesbox}{Exit Ticket --- No Notes}
\small
Bayside Middle School has $750$ students. The principal wants to know whether
students think the lunch period is too short.

\begin{enumerate}[leftmargin=1.6em, itemsep=9pt, topsep=3pt]

  \item Name the main source of bias in each plan. Choose from
        \emph{undercoverage}, \emph{nonresponse}, \emph{social desirability},
        and \emph{demand bias}.

        \begin{center}
        \renewcommand{\arraystretch}{1.45}
        \begin{tabularx}{0.96\linewidth}{@{} X p{3.1cm} @{}}
        \toprule
        \textbf{The plan} & \textbf{Source of bias} \\
        \midrule
        Survey the students who are in the library during lunch.
          & \ans{undercoverage} \\
        Email all $750$ students a link; $88$ of them fill it out.
          & \ans{nonresponse} \\
        A teacher asks students out loud, by name, whether they have ever been
        late to class. & \ans{social desirability} \\
        \bottomrule
        \end{tabularx}
        \end{center}

  \item The principal surveys $30$ students in the library and $24$ of them say
        lunch is too short.

        \textbf{(a)} What percent of those $30$ is that?

        \begin{work}
          \dfrac{24}{30} &= 0.80 \\
                         &= 80\%
        \end{work}

        \textbf{(b)} She repeats it with $60$ library students and gets $48$.
        That is \ans{$80\%$}. Did the larger sample fix the problem? Explain.
        \ansline{No --- it is still only library students, so the same group is missing.}

  \item The survey asks: \emph{``Don't you agree that our rushed $25$-minute
        lunch should be longer?''} Name what is wrong with that wording, and
        rewrite the question so it does not push students either way.
        \ansline{``Don't you agree'' and ``rushed'' tell the student the answer wanted, which invites}
        \ansline{acquiescence. Better: ``Is the 25-minute lunch period long enough? Yes / No.''}

\end{enumerate}
\end{notesbox}

\end{document}
